\chapter*{Glossary}
\markboth{Glossary}{Glossary}
\addcontentsline{toc}{chapter}{Glossary}
\label{revision:glossary}

\begin{description}
  \item[ABI] An application binary interface: the conventions that cooperating
        machine-code routines use for arguments, results, registers, storage and calls.
  \item[Accumulator] State that combines successive contributions, such as a sum
        or checksum. Its width and initial value are part of the algorithm.
  \item[Active window] The selected samples or bytes a call may use. Padding inside
        a containing allocation need not belong to that window.
  \item[Alias] Two references to overlapping storage. An exact alias and a partial
        overlap can have different permissions in a function contract.
  \item[Alignment] A restriction on an address's divisibility. Check the applicable
        interface and instruction rather than assuming every object has one alignment.
  \item[Assembler] A tool that translates assembly source into encoded instructions and data, resolving symbolic addresses and reporting source errors.
  \item[Assembly language] A textual notation for machine instructions, data and assembly directives. Its mnemonics and symbols are source notation, rather than values read directly by the processor.
  \item[Backstep] A simulator operation that restores an earlier recorded execution
        state. It is distinct from running another forward instruction.
  \item[Bit field] A selected sequence of bits within a larger representation.
        Its width and position determine how to extract or replace its value.
  \item[Byte order] The order of bytes representing a multibyte value in memory.
        It is separate from the order in which a format consumes bits within one byte.
  \item[Cache line] A contiguous block transferred and tracked together by a cache. Its size and the mapping policy affect which addresses share or compete for cached storage.
  \item[Caller and callee] The routine making a call and the routine receiving it.
        Their interface divides responsibility for arguments, results and retained state.
  \item[Capacity] The size of available storage. A valid input length or accepted
        output length can be smaller than that capacity.
  \item[Checksum] A compact value computed from input to help detect changes. Equal values alone do not establish that two inputs are identical.
  \item[Clipping] Restricting a mathematical result to specified endpoints. A
        result already equal to an endpoint need not count as clipped.
  \item[Codec] A system for encoding and decoding a representation. Complete
        interoperability involves the stream's syntax and the reconstructed data.
  \item[Continuation] The state or location at which an operation resumes. A
        function's return continuation is one example.
  \item[Descriptor] A handle returned for a resource such as an open file. Its
        numeric value is not the file's contents or its byte position.
  \item[Directive] A source command to the assembler, such as selecting a section or defining data. It is not itself an instruction executed by the simulated processor.
  \item[Endian] Describing byte order: little-endian places a multibyte value's
        least significant byte first; big-endian places its most significant byte first.
  \item[Exclusive end] The boundary just after a selected range. The end itself is
        outside a half-open range.
  \item[Extent] The containing range from an object's start through its last
        selected byte. A strided object's extent can include inactive gaps.
  \item[Fixed point] A representation in which integers have an agreed scale.
        Addition, multiplication, rounding and display must retain that scale's meaning.
  \item[Frame] Context determines the meaning: a function's stack frame is its
        call-related storage, while a video frame is a picture at one sequence position.
  \item[Halfword] A sixteen-bit value or storage unit in the terminology used here.
        Signed and unsigned loads interpret its high bit differently.
  \item[Instruction] An encoded operation defined by the target architecture. Its operands and current machine state determine its architectural effect.
  \item[Invariant] A statement that holds at a chosen point in every iteration or
        state transition and supports the algorithm's correctness argument.
  \item[Label] A source symbol naming a position in code or data. The assembler resolves that position to an address; the name itself occupies no simulated register.
  \item[Lane] One element of a vector register under a selected element width.
        The same register bits can be displayed as different lane widths.
  \item[Mask] A bit pattern used to select or change particular bits. A mask's
        width matters when values occupy wider registers.
  \item[Memory-mapped I/O] Device control or feedback accessed through designated
        memory addresses. A command store and the device's completed action can be separate.
  \item[MSA] The MIPS SIMD Architecture used for vector operations in this course.
        An instruction applies its defined operation to lanes of a vector register.
  \item[Overflow] A mathematical result outside the representable range of an
        operation or storage format. Wrapping, trapping and clipping are different behaviors.
  \item[Prediction] In a codec, a sample estimate used to represent a difference.
        In a learning activity, a reasoned expectation stated before execution.
  \item[Program counter] The register identifying an execution position. A source
        line and an instruction address are related through the assembled program's map.
  \item[Pseudoinstruction] An assembler-supported source operation expanded into one or more actual target instructions. Inspect the assembled form when reasoning about instruction counts or pauses.
  \item[Publication] The point at which prepared results become visible in
        caller-owned output. Failure before that point can require retaining old output.
  \item[Quantization] Selecting a representable level from a wider numerical range.
        The selected level and its reconstructed coefficient have different meanings.
  \item[Register] A named processor storage location read or written by instructions.
        Its current bits do not by themselves specify a signedness or data format.
  \item[Residual] A difference between a current sample and its prediction. Signed
        residuals must retain their sign through the following arithmetic.
  \item[Scalar] Operating on individual values, in contrast with a vector operation
        on a group of lanes. A scalar baseline can help establish vector equivalence.
  \item[Sign extension] Filling a widened integer's new upper bits with copies of
        its original sign bit. Zero extension instead fills those bits with zeros.
  \item[Stack frame] Storage associated with a function call, including any saved
        state, local data and required alignment. Its ownership follows the call contract.
  \item[Stride] The byte advance between corresponding positions in adjacent rows
        or records. It need not equal the count of active bytes in a row.
  \item[Syscall] An instruction that requests a service through the execution environment. In these examples, MARS selects its educational service using the documented register convention.
  \item[Temporal reference] A coded picture's clock-related sequence value. Its
        bounded representation can wrap while the video sequence continues.
  \item[Two's complement] The signed-integer interpretation in which a fixed-width
        negative value is represented by its value modulo that width's power of two.
  \item[Valid length] The amount of meaningful input, distinct from allocated
        capacity. A bitstream can end within its final allocated byte.
  \item[Vector] A register value interpreted as several elements under an operation's
        lane width. Element order and signedness remain part of the operation's contract.
  \item[Wrapping] Retaining the result modulo a fixed power of two. The retained
        bits can differ from a saturated result or from a wider mathematical sum.
\end{description}
